%% kapa: exam=1 round=26 year=2015 period=2 form=A subject=appraisal_law subject_ko=부동산관계법규 seq=4 number=4
%% source: Q-Net 감정평가사 기출문제 (cst00309) artlSeq=5138001 file=extracted/1008/5138001/제26회_감정평가사 1차_기출/제26회_감정평가사 1차_기출/2교시_감평_A형.pdf
%% conversion: kordoc 4.24.51 -> tools/gen1.py (auto); review=proofread-sonnet; flags=none
\begin{question}{4}
국토의 계획 및 이용에 관한 법령상 개발행위허가에 관한 설명으로 옳은 것은? (단, 조례는 고려하지 않음)
\begin{choices}
  \item 도시ㆍ군계획사업에 의한 행위의 경우에도 개발행위허가를 받아야 한다.
  \item 토지의 일부를 공공용지 또는 공용지로 하기 위한 토지의 분할은 개발행위허가를 받아야 한다.
  \item 농림지역에 물건을 1개월 이상 쌓아놓는 경우 개발행위허가를 요하지 아니한다.
  \item 경작을 위한 토지의 형질 변경으로서 전ㆍ답 사이의 지목의 변경을 수반하는 경우에는 개발행위허가를 받아야 한다.
  \item 개발행위허가를 받은 사항으로서 사업면적을 10퍼센트 범위 안에서 축소하는 경우에는 개발행위허가를 요하지 아니한다.
\end{choices}
\end{question}
\answer{3}
